\section{Related Work}
\label{sec:related}

\paragraph{Tool-augmented reasoning and self-correction.}
Tool-augmented language models extend LLMs with access to external tools, APIs, code execution, and symbolic or numerical solvers during reasoning~\citep{schick2023toolformer,yao2022react,lu2023chameleon,qiao2023taskweaver,ferrag2025llm}.
These systems improve performance on tasks requiring computation, retrieval, planning, or environment interaction.
Related work on self-correction and reflection further improves reasoning by asking models to critique, revise, or retry unsuccessful outputs~\citep{renze2024self,madaan2023self,shinn2023reflexion}.
However, tool choice and recovery are often still handled through free-form generation, textual feedback, or iterative retries.

\paragraph{Multi-agent LLM systems and shared memory.}
Multi-agent LLM frameworks decompose complex tasks across role-specialized agents and coordinate them through dialogue, role prompting, message passing, or hand-authored workflows~\citep{wu2024autogen,hong2023metagpt,shen2025exploring,tran2025multi}.
Shared-memory and blackboard-style architectures allow agents to accumulate intermediate results in a common workspace~\citep{han2025exploring,salemi2025llm}.
These systems make intermediate reasoning more accessible, but recovery from intermediate failures is often treated as another round of interaction or as a manually specified fallback.

\paragraph{Learned routing and hierarchical control.}
Selecting among heterogeneous sub-behaviors has long been studied in hierarchical decision-making, options, modular policies, and mixture-of-experts routing~\citep{stolle2002learning,shazeer2017outrageously,habib2025llm,mei2025omnirouter,Yang2025RELOOP}.
These approaches separate high-level control from low-level execution, providing a useful perspective for routing among specialist modules.
In LLM-agent systems, however, routing decisions are often produced by latent generation behavior or task-level prompts rather than by an explicit execution-state transition model.

\paragraph{Spatiotemporal reasoning with language models.}
Recent benchmarks show that LLMs struggle with spatiotemporal reasoning across spatial relations, temporal ordering, forecasting, navigation, trajectories, and graph-structured inference~\citep{quan2025benchmarking,li2025stbench,ni2026streasoner}.
Prior studies have explored temporal reasoning, neuro-symbolic temporal reasoning, multimodal spatiotemporal reasoning, and spatial reasoning with LLMs~\citep{xiong2024large,ge2025tremu,cheng2025v,aghzal2023can, li2026zara}.
These tasks provide a natural testbed for specialist composition because they require heterogeneous computations and expose distinct execution failure modes.

\paragraph{Positioning.}
Our \starname{} lies at the intersection of tool-augmented reasoning, multi-agent coordination, and learned routing.
Its focus is not on adding more specialists or tools, but on making inter-agent control explicit through a failure-aware routing matrix.
By conditioning routing on the current agent, task type, and typed execution status, \starname{} separates the question of what went wrong from the question of which specialist should act next.